\paragraph{Ethical Statement}

We acknowledge that the insights of our work depend on the scope of existing datasets and \dataset, which are by no means exhaustive over the entire potent misuse landscape of LLMs. In particular, the tactics we identified are limited by their source data (\lmsysshort and \wildchatshort), which may not reflect the full spectrum of combinatorial jailbreak tactics employed by real users. Although we mine the jailbreak tactics from real-world data, \dataset contains synthetically composed adversarial prompts generated by \mixtral, \chatgpt, and \gptfour, which do not fully resemble in-the-wild user queries by forms and content and may inherit inherent styles of their source models and filtering heuristics. We encourage future works to explore synthetic data generation more closely aligned with human-written attacks. We also note that while safety training can mitigate many types of risks, certain harmful behaviors are inherently context-dependent and thus can only be guarded against by a layer outside the model itself \cite{ai-safety-not-a-model-property}. 

Finally, we account for the consequences and take appropriate ethical measures of publicly releasing our code and data---the primary purpose and usage of \method and \dataset is to facilitate open resources for model safety enhancement, despite the fact they can elicit harmful responses from models. Thus, we plan to gate the \dataset release behind a content warning and terms agreement limiting usage to researchers who provide a valid justification for their need for \dataset. We believe that the marginal risk~\citet{kapoor2024societal-impact-foundation-models} of releasing \method and \dataset is far outweighed by the benefits of accelerating advances in model safety research. Our work makes a substantial attempt to keep safety research at pace with model capability improvements to mitigate greater risks in the future.

